\documentclass[11pt]{article}
\usepackage{fontspec}
\usepackage[french,english]{babel}
\usepackage{amsmath}
\usepackage{amssymb}
\usepackage{amsthm}
\usepackage{mathrsfs}
\usepackage{tikz-cd}
\usepackage{enumitem}

\newcommand{\printedpage}[1]{\marginpar{\footnotesize p.~#1}}
\newcommand{\msp}[1]{\marginpar{\footnotesize [#1]}}
\newcommand{\bD}{\boldsymbol{\Delta}}
\newcommand{\iddots}{\mathinner{\mkern1mu\raise1pt\hbox{.}\mkern2mu\raise4pt\hbox{.}\mkern2mu\raise7pt\hbox{.}\mkern1mu}}

\begin{document}

\printedpage{4}
% head: A. Grothendieck : Dérivateurs, Ch. VII, Édition des Universités de Montpellier II et Paris VII
% foot: Version 22 décembre 2001

\[
\begin{tikzcd}[column sep=small, row sep=small]
 & & & i \arrow[dl] \arrow[dr] & & & \\
 & & i-1 & & i+1 & & \\
 & \iddots \arrow[dl] & & & & \ddots \arrow[dr] & \\
 i-r & & & & & & i+r, % [sic] i+r -> i+s, s est la longueur de la chaîne contenant u_{i+1}
\end{tikzcd}
\]

where $r$ is the length of the structural chain containing $\underline{u}_i$, $s$ that of the structural chain containing $\underline{u}_{i+1}$. This category is an amalgamated sum $\bD_r \amalg \bD_s$ (glued together along their origin), but is not isomorphic to a $\bD_j$, e.g. for lack of a final object.

This observation proves a). To prove b), it suffices to prove that the total order on $I$ is determined intrinsically, by the choice of \textit{the origin} $s_0$ of $I$ (among these % [sic] ces -> ses, possessif
two endpoints). Suppose that we have already determined intrinsically the vertices $s_0$, $s_1$, $\ldots$, $s_i$ ($i \geq 0$). How to determine $s_{i+1}$? (When $i < n$.) Let $I'$ be the full subcategory of $I$
\msp{page 7}
formed by the $j \in I$ which are not among the $s_j$ \texttt{[rather} $s_k$\texttt{]} already constructed, except the last $s_i$ (thus $s_i \in I'$, $\operatorname{card} I' \geq 2$). Then $I'$ is isomorphic to a standard interval, with $s_i$ corresponding to the origin. One of the categories $i\backslash I$, $I/i$ is one-point, the other isomorphic to a category $\bD_j$, and moreover, $i$ is an initial or final object of $\bD_j$. If it is the initial object $0$, we take for $s_{i+1}$ the immediate successor of $i$ in $\bD_j$. If it is the final object $j$, we take for $s_{i+1}$ the immediate predecessor $j-1$.

If we had started from the target $n$ of $I$, instead of starting from the source, we would have found the opposite total order relation. Assigning the target element the order number $0$, we find a renumbering of $I$, $i$ becoming $n-i$, and the transition arrows which were direct for the initial order relation, become retrograde in the opposite order relation, and conversely. We thus see:

\textbf{Proposition 2.} \textit{If $I$ is a standard interval, there exists a unique pair $(I',\theta)$, where $I'$ is a standard interval, and $\theta$ an isomorphism $I \xrightarrow{\sim} I'$ which takes the target of $I$ into the source of $I'$ (and the source of $I$ into the target of $I'$).}

This interval $I'$ is called the \textit{opposite} interval of $I$. We must see it, intrinsically, as the ``same'' category $I$, % [sic] categorie -> catégorie, accent manquant
but with an opposite ``orientation''.

\msp{page 8}

A clearer way to see this is the following. Let $J$ be an interval category, i.e. isomorphic to a standard interval. Then there are exactly (if $I \not\simeq e$) \textit{two} pairs $(I,\theta)$ of a standard interval $I$ and an isomorphism $\theta : I \xrightarrow{\sim} J$. These two isomorphisms induce on $J$ two total orders, opposite to one another. Each pair is uniquely determined by the said total order, which we shall call % [sic] appelera -> appellera, orthographe
an \textit{orientation} of the interval $J$. Thus $J$ has two orientations, and each determines an \textit{oriented isomorphism} (i.e. respecting
\printedpage{5}
% head: A. Grothendieck : Dérivateurs, Ch. VII, Édition des Universités de Montpellier II et Paris VII
% foot: Version 22 décembre 2001
the orientations) of $J$ with a standard interval (\footnote[2]{\textbf{NB} We shall speak of the \textit{oppositely oriented} interval $I^{\mathrm{opp}}$ of an \textit{oriented} interval $I$, which we must take care not to confuse with the opposite category $I^o$. Here $I^{\mathrm{opp}}$ has the same underlying category as $I$, but with an opposite orientation.}). These two intervals are the same if and only if there exists an automorphism of $J$ (necessarily unique) which interchanges its two orientations, or what amounts to the same thing, which interchanges the two endpoints. (Thus the categories opposite to one another
\[
\begin{tikzcd}[column sep=small, row sep=small]
 & \bullet \arrow[dl] \arrow[dr] & & & & \bullet & \\
 \bullet & & \bullet & & \bullet \arrow[ur] & & \bullet \arrow[ul]
\end{tikzcd}
\]
are each % [sic] chacunes -> chacune, orthographe
symmetric intervals.) Choosing an orientation also amounts to % [sic] `a -> à, accent imprimé comme guillemet
choosing the origin among the two endpoints.

Let us note that the \textit{opposite category} of an interval is an interval, and the opposite category of a standard interval is the standard interval deduced from it by changing all the signs on the transition intervals $[i-1,i]$ ($1 \leq i \leq n$). We shall take care not to confuse the property $I \simeq I^o$, with the property of symmetry.

\msp{page 9}

Thus the two categories ($\Psi_1$ and $\Phi_1$) above are symmetric, but opposite to one another without being isomorphic. On the other hand, the $\bD_i$ are not symmetric, but each isomorphic to the opposite category.

If $I$ is isomorphic to the opposite category $I^o$, and if $\theta$ is such an isomorphism (necessarily unique), then it is immediate that $\theta$ reverses the orientation (otherwise it would be the identity on the vertices, and the transition arrows would be both direct and indirect), hence $\theta(0) = n$, $\theta(1) = n-1$, and we want $\underline{u}_1$ and $\underline{u}_n$ to be of the same sign. But if $I$ is symmetric, $\underline{u}_1$ and $\underline{u}_n$ are of opposite sign. Thus $I$ cannot be both symmetric and isomorphic to the opposite interval \texttt{[rather to the opposite category; exception:} $I = \{0\}$\texttt{]}.

For any category $\bD_i$, this being an interval, there are (if $i > 0$) \textit{two orientations.} As oriented intervals, there are $\bD_i^+$ (with the ordinary orientation, given by the ordinary order on $\operatorname{Ob} \bD_i = [0,i]$), and $\bD_i^-$ with the opposite orientation.

\printedpage{6}
% head: A. Grothendieck : Dérivateurs, Ch. VII, Édition des Universités de Montpellier II et Paris VII
% foot: Version 22 décembre 2001

\msp{page 10}

\section*{2\quad Full subcategories of an interval}

Let $I$ be an oriented interval category, of length $n \geq 1$. We shall identify $\operatorname{Ob} I$ with the set of integers $[0,n]$. For a chain $\underline{u}_i$, $\underline{u}_{i+1}$, $\ldots$, $\underline{u}_{i+l-1}$ of transition arrows in $I$, we call the \textit{domain} of the chain, the set of sources and targets of these elements. This is thus an interval of integers $[i-1, i+l-1]$, of length $l$ equal to that of the chain, whose origin is determined % [sic] déterminé -> déterminée, accord avec l'origine
by the first arrow $\underline{u}_i$ of the chain (it is its source, if $\underline{u}_i$ is direct, and its target if $\underline{u}_i$ is retrograde), and whose endpoint is determined by the last $\underline{u}_{i+l-1}$. The intervals of integers (viewed as vertices of $I$) corresponding to the structural chains are called the \textit{structural intervals} or \textit{critical intervals} of $I$. Unlike the structural chains, these intervals do not form a partition of $I$, since two structural intervals may have a common endpoint. If the structural chains are
\[
C_1,\ C_2,\ \ldots\ ,\ C_r \qquad (r = \text{ useful length of } I)
\]
of lengths
\[
l_1,\ l_2,\ \ldots\ ,\ l_r,
\]
then the structural intervals are
\[
\underbrace{[0,l_1]}_{I_1},\ \underbrace{[l_1,l_1+l_2]}_{I_2},\ \ldots\ ,\ \underbrace{[l_1+\cdots+l_{r-1},l_1+\cdots+l_{r-1}+l_r]}_{I_r},
\]
a sequence of intervals of integers of which two consecutive ones have a common endpoint, and such that if $i < j$, each element of $I$ % [sic] I -> I_i, comparaison des éléments de I_i et de I_j
is $\leq$ each element of $I_j$. We call \textit{critical vertex} of $I$ the endpoints of the structural intervals.

\textbf{Proposition 3.} \textit{Let $s \in I$, let $I' = I/s$ (full subcategory of $I$ of the $s'$ such that there exists $s' \longrightarrow s$), $I'' = I/s$ \textup{\texttt{[rather} $I'' = s\backslash I$\texttt{]}}, $I'_0 = (I/s)\backslash\{s\}$, $I''_0 = (s\backslash I)\backslash\{s\}$} (\footnote[3]{i.e. that we have: there are no arrows between an element of $I'_0$ and an element of $I''_0$.}). For $s$ to be critical, it is necessary and sufficient that $I\backslash\{s\}$ be the sum category of $I'_0$ and $I''_0$.

\msp{page 11}

\textbf{Proposition 4.} \textit{Let $i,j$ be in $I$, with $i < j$. Equivalent conditions:}
\begin{enumerate}
\item[\textit{a)}] \textit{There exists an arrow between $i$ and $j$ (either $i \longrightarrow j$, or $j \longrightarrow i$).}
\item[\textit{b)}] \textit{$i$ and $j$ belong to the same structural interval (necessarily unique).}
\item[\textit{c)}] \textit{The full subcategory of $I$ defined by the set of objects $[i,j]$ is isomorphic to a $\bD_\alpha$ (we shall have $\alpha = j-i$).}
\end{enumerate}

\printedpage{7}
% head: A. Grothendieck : Dérivateurs, Ch. VII, Édition des Universités de Montpellier II et Paris VII
% foot: Version 22 décembre 2001
We shall then say that $i$ and $j$ are \textit{directly connected} in $I$. (And we shall use the same terminology in an arbitrary category: $i$ and $j$ are directly connected if and only if $\operatorname{Hom}_I(i,j)$ \textit{or} $\operatorname{Hom}_I(j,i)$ is nonempty.)

\textbf{Proposition 5.} \textit{Let $J \subseteq I$ be a subset of $I$, let us consider the full subcategory of $I$ which it defines. Then:}
\begin{enumerate}
\item[\textit{a)}] \textit{The connected components $J_\alpha$ of $J$ are intervals and the order induced on a $J_\alpha$ by that of $I$ is an orientation (called the \textup{induced orientation}) of $J_\alpha$.}
\item[\textit{b)}] \textit{Suppose $\operatorname{card} J > 1$. Let $i$ be the smallest element of $J$, $j$ the largest element, thus $i < j$. For $J$ itself to be an interval, it is necessary and sufficient that every critical vertex in $[i,j]$ be an element of $J$.}
\end{enumerate}

(\footnote[4]{We call a \textit{subinterval} of $I$ a full subcategory % [sic] sous-categorie -> sous-catégorie, accent manquant
which is connected, i.e. which is itself an interval. Let us note moreover that \textit{every} subcategory $I'$ of $I$ which is an interval, is necessarily full, cf. below \texttt{[paragraph]} on functors between intervals.}).

\textbf{NB} If $s \in [i,j]$ is a critical vertex of $I$ which is not in $J$, then by proposition 3 it would disconnect $J$, i.e. $J$ would be a direct sum of two nonempty subcategories $J'$ and $J''$ ($i \in J'$, $j \in J''$), hence $J$ would not be an interval.

\msp{page 12}

\textbf{Corollary.} \textit{The set of critical vertices of $I$ is also the set of vertices of the smallest subinterval of $I$ containing the endpoints. More generally, if $i < j$ are two vertices of $I$, the smallest subinterval of $I$ containing $i$, $j$ is formed by $i$, $j$ and the critical vertices $s$ such that $i < s < j$.}

\textbf{Definition.} An interval is said to be \textit{reduced} (or \textit{irredundant}) if all its vertices are critical, or what amounts to the same thing, if its structural chains are of length 1.

A \textit{reduced} oriented interval is known, up to isomorphism, when we know
\begin{enumerate}
\item[a)] its length $\mathrm{n}$,
\item[b)] its initial direction $+$ (direct) or $-$ (retrograde).
\end{enumerate}

\textbf{NB} This direction is the same as the final direction if $n$ is odd, it is opposite if $n$ is even.

Thus, according to the two cases of initial direction, we find
\[
\begin{aligned}
0 \longrightarrow 1 \longleftarrow 2 \longrightarrow \cdots \qquad\qquad \cdots\ n & \qquad (\text{With } & n-1 \longrightarrow n & \quad \text{if } n \text{ odd,}\\
& & n-1 \longleftarrow n & \quad \text{if } n \text{ even.)}\\
0 \longleftarrow 1 \longrightarrow 2 \longleftarrow \cdots \qquad\qquad \cdots\ n & \qquad (\text{With } & n-1 \longleftarrow n & \quad \text{if } n \text{ odd,}\\
& & n-1 \longrightarrow n & \quad \text{if } n \text{ even.)}
\end{aligned}
\]

\printedpage{8}
% head: A. Grothendieck : Dérivateurs, Ch. VII, Édition des Universités de Montpellier II et Paris VII
% foot: Version 22 décembre 2001

\textbf{Proposition 6.} \textit{A subinterval $J$, having $i$ and $j$ ($i < j$) as smallest and largest vertex, is reduced if and only if every $s \in J$ distinct from $i$ and from $j$ is critical.}

\textbf{Corollary.} \textit{The full subcategory of $I$ formed by the critical vertices is a subinterval. We call it the reduced interval associated to $I$, denoted $I_{\textup{red}}$.}

\end{document}
